\chapter{AEO --- Answer Engine Optimisation}
\label{ch:aeo}

\section{Applicability}

Answer Engine Optimisation concerns how readily a product's content is selected
and quoted by answer engines --- AI assistants, conversational search and
featured-snippet surfaces. As with SEO (Chapter~\ref{ch:seo}), there is no
project-owned website, so the question reduces to: \emph{how well does the
repository and its listings feed answer engines?}

\section{Content available to answer engines}

Answer engines draw on: the README files, the package metadata, the recipe
descriptions, the extension CHANGELOG, and any prose in the repository. All of
these are plain text or Markdown, which is exactly the format answer engines can
consume and quote.

\begin{longtable}{p{4.4cm}p{3.2cm}p{6.6cm}}
\caption{Answer-engine-relevant content}\\
\toprule
\textbf{Source} & \textbf{Format} & \textbf{Value for answer engines} \\
\midrule

Extension README & Markdown & Explains what the tool does, the problem it solves,
and how to use it, in natural language. \\
\texttt{ProjectDescription.md} & Markdown & A product description with framing that
answers ``what is this?''. \\
Recipe \texttt{description} fields & Plain text & One-line answers to ``is there a
screen for X?''. \\
Extension CHANGELOG & Markdown & Dated, factual statements suitable for quoting. \\
Package descriptions & Plain text & Short capability statements. \\
\bottomrule
\end{longtable}

\section{Answer-readiness of the content}

\begin{longtable}{p{1.5cm}p{4.6cm}p{3.4cm}p{4.7cm}}
\caption{AEO findings}\\
\toprule
\textbf{ID} & \textbf{Finding} & \textbf{Evidence} & \textbf{Recommendation} \\
\midrule

AEO-01 & The repository contains clear ``what it is / what problem'' prose in the
extension README, which is well suited to being quoted verbatim. & Extension
README sections ``What is Expo Snips?'' and ``The Problem''. & IMPLEMENTED; keep
this structure. \\
AEO-02 & FAQ-shaped content is absent. Answer engines favour explicit question-and-answer
pairings that can be lifted directly. & No FAQ section anywhere. & RECOMMENDED:
add a short FAQ (``Does it work offline?'', ``Which Expo SDK?'', ``Is it free?''). \\
AEO-03 & \texttt{documents/ProjectDescription.md} contains \emph{outdated}
statements: it references a recipe count and variant support that the code no
longer matches. & \texttt{documents/ProjectDescription.md}. & Update it, because
an answer engine may quote the stale figure. See
Chapter~\ref{ch:known-issues}. \\
AEO-04 & Authoritative one-sentence definitions are duplicated in several places
(README, package descriptions) with slightly different wording, which can produce
inconsistent answers. & \texttt{vscode-extension/package.json} description vs
README headline. & Standardise a single canonical one-liner. \\
AEO-05 & No structured data (schema.org \texttt{SoftwareApplication} or
\texttt{SoftwareSourceCode}) is published, because there is no site to host it. &
Absence of any HTML. & NOT APPLICABLE unless a site is created; then add JSON-LD. \\
AEO-06 & The README uses license and version badges that are useful, unambiguous
facts for an answer engine to extract. & Badge block in the extension README. &
IMPLEMENTED. \\
\bottomrule
\end{longtable}

\section{Factual accuracy as an AEO risk}

The single most important AEO property is accuracy; an answer engine that quotes
a wrong fact damages trust more than one that quotes nothing. Two facts about
this repository are stated inconsistently across sources:

\begin{itemize}
  \item \textbf{Recipe count:} \texttt{ProjectDescription.md} understates the
        catalog relative to the code, where the count is 24.
  \item \textbf{Variant support:} \texttt{ProjectDescription.md} implies variants
        are unsupported, while three recipes (\texttt{app/home},
        \texttt{auth/login}, \texttt{onboarding/welcome}) implement a
        \texttt{glass} variant.
\end{itemize}

\textbf{Status:} INFERRED from comparison of the document with the code.
\textbf{Recommendation:} correct the description document (RECOMMENDED).

\section{Summary}

\begin{table}[H]
\centering
\caption{AEO posture}
\begin{tabular}{p{4.6cm}p{9.6cm}}
\toprule
\textbf{Dimension} & \textbf{Assessment} \\
\midrule
Quotable capability prose & IMPLEMENTED (extension README). \\
FAQ content & ABSENT (recommended). \\
Structured data & NOT APPLICABLE (no site). \\
Factual consistency & GAP (outdated description document). \\
\bottomrule
\end{tabular}
\end{table}

\textbf{Status:} PARTIAL / RECOMMENDED.
